\chapter{A Concrete Naming Certificate for $\PaperLeanDriftWitnessUp$}
\label{ch:concrete-instances-paper-lean-drift-witness-namecert}
\origin{human}

The $\PaperLeanDriftWitnessUp$ packet names the finite paper-Lean resolution surface used by the drift audit. It records a paper marker occurrence, the normalized declaration name requested by that marker, the duplicate-label refusal row, the Lean inventory lookup row, the resolution verdict, replay provenance, and the local $\NameCert$ row in one certificate. The packet makes paper-marker resolution a concrete carrier rather than a host-side impression that a nearby Lean declaration exists. Its single-site exact-resolution verdict is a local row consumed by the audit-map dependency routing surface of \autoref{ch:concrete-instances-auditmapdependencyweave-namecert}.

\begin{definition}[Paper-Lean drift witness carrier]
\label{def:paper-lean-drift-witness-carrier}
A $\PaperLeanDriftWitnessUp$ carrier is a finite $\BHist$ packet
$$
W=(M,A,L,I,R,H,C,P,N)
$$
with the following displayed rows:
\begin{enumerate}
\item $M$ is the paper marker row, recording the exact manuscript site, marker kind, and source spelling of a \texttt{\textbackslash leanchecked}, \texttt{\textbackslash leanvariant}, \texttt{\textbackslash leanstmt}, \texttt{\textbackslash leandef}, or \texttt{\textbackslash leantarget} occurrence;
\item $A$ is the normalized requested name row obtained from $M$, including underscore normalization and the fully qualified Lean declaration path requested by the paper site;
\item $L$ is the paper-label ledger row, recording the local label attached to the manuscript site and the duplicate-label refusal verdict for the full paper label inventory;
\item $I$ is the Lean inventory lookup row, recording the finite declaration table against which $A$ is checked;
\item $R$ is the resolution verdict row, recording exact resolution, unresolved-name refusal, ambiguous-name refusal, or duplicate-label refusal for the same marker site;
\item $H$ records componentwise $\hsame$ transport for markers, normalized names, label ledgers, inventory rows, verdicts, replay routes, provenance, and local names;
\item $C$ records the displayed $\Cont$ routes by which a consumer replays marker parsing, name normalization, duplicate-label comparison, inventory lookup, and verdict exposure;
\item $P$ is the $\Pkg$ provenance over paper markers, label ledgers, Lean declaration inventory, audit replay, $\hsame$, $\Cont$, $\Pkg$, and $\NameCert$ source packets;
\item $N$ is the local $\NameCert_{\PaperLeanDriftWitnessUp}$ row.
\end{enumerate}
The carrier has no coordinate for approximate name matching, silent marker repair, duplicate-label acceptance, unresolved-name promotion, or a checked-looking paper claim without an exact Lean inventory hit.
\end{definition}

\begin{theorem}[Paper-Lean drift witness NameCert obligations]
\label{thm:paper-lean-drift-witness-namecert-obligations}
For every accepted $\PaperLeanDriftWitnessUp$ carrier
$$
W=(M,A,L,I,R,H,C,P,N),
$$
the local $\NameCert_{\PaperLeanDriftWitnessUp}$ surface consists exactly of marker-site exposure, normalized requested-name exposure, duplicate-label refusal, Lean inventory lookup, resolution verdict exposure, componentwise $\hsame$ transport, displayed $\Cont$ replay, $\Pkg$ provenance, and the local naming row. It does not export a separate paper authority for treating a non-resolving marker as checked.
\end{theorem}
\leanchecked{BEDC.Derived.PaperLeanDriftWitnessUp.PaperLeanDriftWitness\_namecert\_obligations}
\begin{proof}
Carrier exposure is projection from \autoref{def:paper-lean-drift-witness-carrier}. The tuple lists the paper marker row, normalized name row, label ledger, Lean inventory lookup, verdict row, transport, replay, provenance, and local name.

The NameCert obligations are the finite audit obligations on those rows. The marker and name rows expose what the paper asked for; the label ledger and inventory row expose the two finite tables against which the request is checked; the verdict row records the admissible resolution outcome.

Transport and replay are confined to $H$ and $C$, with provenance and naming confined to $P$ and $N$. Since no row stores approximate matching, silent repair, duplicate-label acceptance, or unresolved-name promotion, the local certificate is exhausted by the displayed obligations.
\end{proof}

\begin{theorem}[Paper-Lean drift witness resolution exactness]
\label{thm:paper-lean-drift-witness-resolution-exactness}
If an accepted $\PaperLeanDriftWitnessUp$ packet records an exact-resolution verdict in $R$, then the normalized requested name $A$ is present as the same fully qualified declaration in the Lean inventory row $I$, and the paper label row $L$ is not refused as a duplicate label. The packet never identifies a marker with a neighboring declaration, a spelling variant not present in $I$, or a label site that the duplicate-label ledger refuses.
\end{theorem}
\leanchecked{BEDC.Derived.PaperLeanDriftWitnessUp.PaperLeanDriftWitness\_resolution\_exactness}
\begin{proof}
Read the accepted packet through the carrier. The exact-resolution verdict is a row of $R$ attached to the same marker $M$, normalized name $A$, label ledger $L$, and inventory row $I$.

The replay route $C$ first normalizes the marker spelling, then compares the label row against the paper ledger, and then looks up the normalized name in the Lean declaration inventory. Exact resolution can therefore be read only after both finite comparisons have admitted the site.

Componentwise transport in $H$ preserves the same marker, label, and inventory rows. Thus a consumer cannot replace the resolved name by a neighboring declaration, a spelling not present in $I$, or a label site that $L$ refuses.
\end{proof}

\begin{theorem}[Paper-Lean drift witness non-escape]
\label{thm:paper-lean-drift-witness-nonescape}
An accepted $\PaperLeanDriftWitnessUp$ packet exports no checked theorem, statement, definition, or closure-status target beyond the exact marker-resolution rows
$$
M,A,L,I,R,H,C,P,N.
$$
In particular, it exports no unresolved Lean target, no duplicate paper label, no marker-kind coercion, no approximate declaration match, and no host-side permission to treat paper prose as a resolving Lean declaration.
\end{theorem}
\begin{proof}
Inspect the carrier inventory. The rows record marker parsing, requested-name normalization, label-ledger comparison, Lean inventory lookup, resolution verdict, transport, replay, provenance, and local naming.

None of these rows is a theorem proof, a definition body, a closure-status upgrade, or a host permission row. The positive content is only the exact resolution verdict attached to a marker whose normalized name occurs in the inventory and whose label row is not refused.

Therefore any unresolved target, duplicate label, marker-kind coercion, approximate match, or prose-only declaration leaves the packet. Such content must be supplied by another explicit certificate and cannot be projected from this drift witness.
\end{proof}

\begin{theorem}[Paper-Lean drift witness consumer boundary]
\label{thm:paper-lean-drift-witness-consumer-boundary}
Every drift-audit consumer of an accepted $\PaperLeanDriftWitnessUp$ packet factors through the paper marker row $M$, normalized requested name $A$, duplicate-label ledger $L$, Lean inventory lookup $I$, resolution verdict $R$, componentwise transport $H$, replay route $C$, provenance row $P$, and local $\NameCert$ row $N$. The consumer receives a finite NameCert surface for exact marker resolution, not a general semantic equivalence between paper claims and Lean declarations.
\end{theorem}
\begin{proof}
Start with the carrier. The audit consumer can project only the nine displayed rows $M,A,L,I,R,H,C,P,N$.

Resolution reads pass through the normalized name and Lean inventory rows. Duplicate-label reads pass through the label ledger. Verdict reads pass through $R$ and can be replayed only by the displayed routes in $C$ under componentwise transport $H$.

The provenance row $P$ and local name row $N$ package this finite read as a NameCert surface. Since the carrier contains no semantic-equivalence row between paper prose and Lean theorem meaning, the consumer obtains exact marker resolution only.
\end{proof}

\begin{theorem}[Paper-Lean drift witness verdict determinism]
\label{thm:paper-lean-drift-witness-verdict-determinism}
For an accepted $\PaperLeanDriftWitnessUp$ carrier
$$
W=(M,A,L,I,R,H,C,P,N),
$$
fix the paper marker row $M$, normalized requested-name row $A$, paper-label ledger row $L$, and Lean inventory row $I$. The verdict row $R$ exposes exactly one local outcome for that same quadruple: exact resolution, unresolved-name refusal, ambiguous-name refusal, or duplicate-label refusal. A consumer replaying the displayed route $C$ under componentwise transport $H$ cannot obtain a second verdict for the same $M,A,L,I$.
\end{theorem}
\begin{proof}
Project the accepted carrier to the rows $M,A,L,I,R,H,C,P,N$. The rows $M,A,L,I$ determine the marker site, the normalized request, the paper-label ledger read, and the finite Lean inventory read being adjudicated.

The verdict coordinate is a single row $R$ attached to those projected inputs. Its allowed local outcomes are the four carrier outcomes listed in \autoref{def:paper-lean-drift-witness-carrier}: exact resolution, unresolved-name refusal, ambiguous-name refusal, and duplicate-label refusal.

Replay does not introduce another decision channel. The displayed route $C$ parses the marker, normalizes the requested name, checks the label ledger, checks the Lean inventory, and exposes the verdict row. Componentwise transport in $H$ preserves the same inputs through that route, so replaying the same carrier cannot expose a distinct verdict for the fixed $M,A,L,I$.
\end{proof}
\leanchecked{BEDC.Derived.PaperLeanDriftWitnessUp.PaperLeanDriftWitness\_verdict\_determinism}

\begin{theorem}[Paper-Lean drift witness duplicate-label priority]
\label{thm:paper-lean-drift-witness-duplicate-label-priority}
Let
$$
W=(M,A,L,I,R,H,C,P,N)
$$
be an accepted $\PaperLeanDriftWitnessUp$ carrier. If the paper-label ledger row $L$ records duplicate-label refusal for the marker site $M$, then no exact-resolution verdict in $R$ may be exported for that site, even when the normalized requested name $A$ is present in the Lean inventory row $I$.
\end{theorem}
\begin{proof}
Read the carrier in the displayed row order $M,A,L,I,R,H,C,P,N$. The label ledger row $L$ is consulted before the verdict row is exposed, and its duplicate-label refusal is a local refusal for the same marker site $M$.

By \autoref{thm:paper-lean-drift-witness-resolution-exactness}, exact resolution requires both an inventory hit for $A$ in $I$ and the absence of duplicate-label refusal in $L$. Thus an inventory hit alone is not sufficient for exporting an exact-resolution row.

The replay route $C$ keeps the label-ledger comparison on the path to $R$, and $H$ preserves that ledger row under transport. Therefore a duplicate-label refusal in $L$ has priority over the inventory hit and prevents export of an exact-resolution verdict for the refused marker site.
\end{proof}

\begin{theorem}[Paper-Lean drift witness unresolved marker refusal]
\label{thm:paper-lean-drift-witness-unresolved-marker-refusal}
Let
$$
W=(M,A,L,I,R,H,C,P,N)
$$
be an accepted $\PaperLeanDriftWitnessUp$ carrier. If the normalized requested name $A$ is absent from the Lean inventory row $I$, then every consumer read through $R,H,C,P,N$ remains a refusal row and cannot be promoted to a checked marker.
\end{theorem}
\begin{proof}
Absence of $A$ from $I$ blocks the exact-resolution condition. By \autoref{thm:paper-lean-drift-witness-resolution-exactness}, an exact-resolution verdict can be read only when the normalized requested name is present in the Lean inventory row and the paper-label ledger does not refuse the site.

The only remaining verdict information for the same marker passes through the refusal rows of $R$. The route $C$ replays the same normalization and inventory lookup, while $H$ preserves the same normalized name and inventory row under transport.

Non-escape from \autoref{thm:paper-lean-drift-witness-nonescape} prevents the packet from exporting a checked theorem, statement, definition, or closure-status target outside the exact marker-resolution rows. The consumer boundary in \autoref{thm:paper-lean-drift-witness-consumer-boundary} confines downstream reads to $R,H,C,P,N$, so the absent inventory hit remains a refusal and cannot become a checked marker.
\end{proof}

\begin{theorem}[Paper-Lean drift witness resolution stability]
\label{thm:paper-lean-drift-witness-resolution-stability}
Let
$$
W=(M,A,L,I,R,H,C,P,N)
$$
be an accepted $\PaperLeanDriftWitnessUp$ carrier with an exact-resolution verdict in $R$. If
$$
M \hsame M',\quad A \hsame A',\quad L \hsame L',\quad I \hsame I',\quad R \hsame R',\quad H \hsame H',\quad C \hsame C',\quad P \hsame P',\quad N \hsame N',
$$
are componentwise transport rows displayed by $H$, then the transported packet
$$
W'=(M',A',L',I',R',H',C',P',N')
$$
has the same exact-resolution outcome for the transported marker site. In particular, exact resolution remains an inventory hit for $A'$ in $I'$ together with the absence of duplicate-label refusal in $L'$.
\end{theorem}
\begin{proof}
The carrier theorem exposes the nine coordinates as the complete local surface for a drift witness. The NameCert-obligation theorem identifies $H$ as componentwise transport for marker rows, normalized names, label ledgers, inventory rows, verdicts, replay routes, provenance, and local names.

Apply the resolution-exactness theorem before transport. Exact resolution of $W$ reads the same marker through $M$, the normalized request $A$, the duplicate-label ledger $L$, and the Lean inventory row $I$; it requires an inventory hit and no duplicate-label refusal.

The transport rows in $H$ carry each of these coordinates to the corresponding primed row. Non-escape prevents an additional row from changing the verdict during transport, and the consumer-boundary theorem forces every downstream read of $W'$ through the transported coordinates.

Therefore the transported verdict in $R'$ is still exact resolution for $M'$ precisely by the transported inventory hit for $A'$ in $I'$ and the transported non-refusal in $L'$.
\end{proof}

\begin{theorem}[Paper-Lean drift witness marker-kind separation]
\label{thm:paper-lean-drift-witness-marker-kind-separation}
For every accepted $\PaperLeanDriftWitnessUp$ carrier
$$
W=(M,A,L,I,R,H,C,P,N),
$$
the marker kind recorded by $M$ is preserved as its displayed paper-row kind throughout normalization, inventory lookup, verdict exposure, transport, replay, provenance, and local naming. A row whose source marker is \texttt{\textbackslash leanchecked} cannot be consumed as \texttt{\textbackslash leanvariant}, \texttt{\textbackslash leanstmt}, \texttt{\textbackslash leandef}, or \texttt{\textbackslash leantarget}; the same separation holds for each of the five marker kinds.
\end{theorem}
\begin{proof}
The carrier definition stores the marker kind inside the paper marker row $M$. The NameCert-obligation theorem exposes this marker-site row as one of the required certificate rows and gives no separate coordinate for marker-kind conversion.

Resolution exactness then reads the normalized declaration request from $A$ as obtained from that same marker row, and it checks $A$ against the inventory row $I$ under the duplicate-label ledger $L$. The verdict is therefore attached to the marker site and kind already exposed by $M$.

Non-escape rules out marker-kind coercion as exported content of the packet. The replay and transport coordinates $C$ and $H$ can only replay and transport the displayed rows, while the consumer-boundary theorem confines consumers to $M,A,L,I,R,H,C,P,N$.

Thus no consumer can reinterpret a row for one of \texttt{\textbackslash leanchecked}, \texttt{\textbackslash leanvariant}, \texttt{\textbackslash leanstmt}, \texttt{\textbackslash leandef}, or \texttt{\textbackslash leantarget} as a row for another kind.
\end{proof}

\begin{theorem}[Paper-Lean drift witness inventory verdict totality]
\label{thm:paper-lean-drift-witness-inventory-verdict-totality}
For every accepted $\PaperLeanDriftWitnessUp$ carrier
$$
W=(M,A,L,I,R,H,C,P,N),
$$
each accepted verdict exposed by $R$ is read from the displayed Lean inventory row $I$ together with the duplicate-label ledger row $L$. No accepted verdict may be obtained from approximate declaration matching, prose-only promotion, or a paper marker whose duplicate-label row is refused.
\end{theorem}
\begin{proof}
The carrier definition lists the admissible resolution outcomes in the verdict row $R$ and lists no coordinate for approximate matching or prose-only promotion. The NameCert obligations require the duplicate-label ledger, Lean inventory lookup, and verdict exposure as displayed rows of the same packet.

For the exact-resolution case, apply the resolution-exactness theorem. The verdict requires that $A$ occurs as the same fully qualified declaration in $I$ and that $L$ does not refuse the marker site.

For refusal cases, the non-escape theorem prevents unresolved names, duplicate labels, marker-kind coercions, approximate matches, and prose-only declarations from leaving the exact marker-resolution surface as checked content. The consumer-boundary theorem confines every audit-map consumer to the rows $M,A,L,I,R,H,C,P,N$.

Hence every accepted verdict is total over the displayed inventory and duplicate-label ledger rows: it is either the exact inventory-and-ledger outcome or a refusal exposed by $R$, never an approximate or prose-only promotion.
\end{proof}

\begin{theorem}[Paper-Lean drift witness token objectwise audit]
\label{thm:paper-lean-drift-witness-token-objectwise-audit}
For every accepted $\PaperLeanDriftWitnessUp$ packet
$$
W=(M,A,L,I,R,H,C,P,N),
$$
any consumer using $W$ to justify a \texttt{\textbackslash formalstatus\{...\}} token or a marker-resolution synchronization must read the exact marker row $M$, normalized requested-name row $A$, duplicate-label ledger row $L$, Lean inventory row $I$, and verdict row $R$ for that same paper site. The packet cannot convert objective Lean inventory strength into an aggregate route-level formalstatus, bridge row, theoryclosure row, or semantic equivalence between paper prose and Lean declarations.
\end{theorem}
\begin{proof}
The positive resolution case is objectwise. By resolution exactness, an exact inventory hit is attached to the same marker row $M$, normalized name row $A$, label ledger row $L$, and Lean inventory row $I$, and it is exposed only as the verdict row $R$ for that site.

Verdict determinism and inventory verdict totality then fix the local outcome. For the fixed rows $M,A,L,I$, the consumer can read one displayed verdict from $R$: either the exact inventory-and-ledger outcome or a refusal row over the same local inventory and duplicate-label ledger.

The consumer boundary confines downstream reads to $M,A,L,I,R,H,C,P,N$. Non-escape excludes checked theorem content, closure-status targets, bridge rows, theoryclosure rows, and semantic-equivalence content not present in those rows.

Therefore a \texttt{\textbackslash formalstatus\{...\}} token or marker-resolution synchronization justified by $W$ is evidence about the displayed objectwise marker-resolution rows for one paper site. It is not aggregate route-level closure evidence and cannot collapse the manuscript closure axis with the Lean inventory lookup axis.
\end{proof}

\begin{theorem}[Paper-Lean drift witness audit-map dependency source lock]
\label{thm:paper-lean-drift-witness-auditmap-dependency-source-lock}
Let
$$
W=(M,A,L,I,R,H,C,P,N)
$$
be an accepted $\PaperLeanDriftWitnessUp$ carrier. Any $\AuditMapDependencyWeaveUp$ consumer using $W$ as a dependency source must read the same marker row $M$, normalized requested-name row $A$, duplicate-label ledger row $L$, Lean inventory row $I$, verdict row $R$, componentwise transport row $H$, replay row $C$, provenance row $P$, and local $\NameCert$ row $N$. The consumer cannot replace those rows by approximate marker repair, an aggregate formalstatus token, a neighboring Lean name, or a route-level closure summary.
\end{theorem}
\begin{proof}
The carrier projection exposes exactly the rows $M,A,L,I,R,H,C,P,N$. The marker row, normalized request, ledger, inventory, and verdict are therefore read from the same accepted witness rather than from a neighboring audit route.

Resolution exactness fixes the positive read. An exact verdict requires the displayed normalized name $A$ to occur in the displayed inventory $I$ and requires the displayed ledger $L$ not to refuse the same marker site $M$.

Token objectwise audit prevents this local verdict from becoming an aggregate formalstatus route. The consumer boundary then confines every downstream $\AuditMapDependencyWeaveUp$ read to the same displayed rows, so approximate marker repair and route-level closure summaries are outside the handoff.
\end{proof}

\begin{theorem}[Paper-Lean drift witness objectwise drift handoff]
\label{thm:paper-lean-drift-witness-objectwise-drift-handoff}
For every accepted $\PaperLeanDriftWitnessUp$ carrier, an exact verdict hands off only the objectwise marker row $M$, normalized Lean name row $A$, duplicate-label ledger row $L$, Lean inventory hit row $I$, and verdict row $R$ for the same paper site. Replay, transport, provenance, and local naming may carry that row through $H,C,P,N$, but they cannot add a semantic-equivalence row, an approximate declaration match, or a checked-looking marker outside the exact verdict.
\end{theorem}
\begin{proof}
Resolution exactness supplies the objectwise positive case: the marker row $M$, normalized request $A$, ledger $L$, inventory $I$, and verdict $R$ are read together for one paper site.

Duplicate-label priority keeps the ledger in the handoff. If $L$ refuses the site, an inventory-shaped hit cannot be exported as an exact verdict, so the consumer cannot ignore the paper-label row.

Unresolved-marker refusal keeps absent inventory hits inside refusal rows. Verdict determinism and inventory verdict totality then make the exported row exhaustive over the displayed local outcomes, leaving only transport, replay, provenance, and naming through $H,C,P,N$.
\end{proof}

\begin{theorem}[Paper-Lean drift witness audit-map weave consumer exactness]
\label{thm:paper-lean-drift-witness-auditmap-weave-consumer-exactness}
For every accepted $\PaperLeanDriftWitnessUp$ carrier, the downstream $\AuditMapDependencyWeaveUp$ surface receives an exact local dependency row for marker resolution: the source-locked packet $M,A,L,I,R,H,C,P,N$ and the objectwise verdict handoff over $M,A,L,I,R$. It receives no semantic equivalence between paper prose and Lean declarations, no checked proof of the named Lean declaration, and no bridge or closure-axis certificate.
\end{theorem}
\begin{proof}
First apply the audit-map dependency source lock. The dependency source is the displayed finite packet $M,A,L,I,R,H,C,P,N$, and every consumer read must factor through those rows.

Next apply the objectwise drift handoff. The exact verdict exported to the dependency weave is the local marker-resolution row over $M,A,L,I,R$, carried only by the displayed transport, replay, provenance, and local naming rows.

Finally use the consumer-boundary and non-escape theorems. The boundary gives exact marker resolution as dependency data, while non-escape prevents semantic equivalence, Lean proof content, bridge checking, or closure-axis certification from leaving the $\PaperLeanDriftWitnessUp$ packet.
\end{proof}

\begin{theorem}[Paper-Lean drift witness obligation-closure package]
\label{thm:paper-lean-drift-witness-obligation-closure-package}
For every accepted $\PaperLeanDriftWitnessUp$ carrier
$$
W=(M,A,L,I,R,H,C,P,N),
$$
the obligation-closure package is exactly the conjunction of marker-site exposure, normalized requested-name exposure, duplicate-label priority, unresolved-marker refusal, marker-kind separation, inventory verdict totality, componentwise $\hsame$ transport, displayed $\Cont$ replay, $\Pkg$ provenance, and local $\NameCert$ naming. It exports no checked theorem content, no semantic equivalence between paper claims and Lean declarations, no approximate declaration match, and no prose-only promotion.
\end{theorem}
\begin{proof}
The carrier definition displays the nine rows $M,A,L,I,R,H,C,P,N$. The NameCert-obligation theorem identifies those rows as the marker site, normalized requested name, paper-label ledger, Lean inventory lookup, verdict exposure, componentwise transport, replay, provenance, and local naming surface.

Resolution exactness supplies the exact inventory-and-ledger read for an accepted exact verdict. Duplicate-label priority blocks an exact-resolution export when the paper-label ledger refuses the site, while unresolved-marker refusal keeps an absent Lean inventory hit inside a refusal row.

Marker-kind separation prevents a consumer from retyping one paper marker kind as another. Inventory verdict totality then exhausts the accepted verdicts over the displayed Lean inventory row and duplicate-label ledger, ruling out approximate declaration matching and prose-only promotion.

It remains only to bind the package boundary. Componentwise $\hsame$ transport is confined to $H$, displayed $\Cont$ replay is confined to $C$, $\Pkg$ provenance is confined to $P$, and local naming is confined to $N$. By non-escape and the consumer-boundary theorem, no checked theorem content, semantic equivalence row, approximate match, or prose-only promotion can leave this finite witness packet.
\end{proof}

\begin{theorem}[Paper-Lean drift witness closure-target resolution]
\label{thm:paper-lean-drift-witness-closure-target-resolution}
For every accepted $\PaperLeanDriftWitnessUp$ carrier
$$
W=(M,A,L,I,R,H,C,P,N),
$$
any closure-status consumer or paper-marker consumer must read the exact paper marker row $M$, normalized Lean target row $A$, duplicate-label ledger row $L$, Lean inventory row $I$, and verdict row $R$ before treating a \texttt{\textbackslash leanchecked}, \texttt{\textbackslash leanvariant}, \texttt{\textbackslash leanstmt}, \texttt{\textbackslash leandef}, or \texttt{\textbackslash leantarget} occurrence as resolving. The resolving read is accepted only as the displayed exact inventory-and-ledger verdict for that same marker site.
\end{theorem}
\begin{proof}
Begin with the carrier row order. The marker site is exposed by $M$, the requested Lean declaration is normalized in $A$, the paper-label ledger is read through $L$, the Lean inventory is read through $I$, and the verdict is exposed by $R$.

Resolution exactness supplies the positive case: an exact verdict requires the normalized requested name to occur in $I$ and the label row not to be refused by $L$. Duplicate-label priority prevents an inventory hit from resolving a marker whose label row is refused.

Unresolved-marker refusal keeps absent inventory hits inside refusal rows, while marker-kind separation keeps each of the five marker commands in its own paper-row kind. Inventory verdict totality then exhausts the accepted outcomes over the displayed inventory and ledger rows.

Therefore a closure-status or marker consumer can treat a paper marker as resolving only after reading $M,A,L,I,R$ for the same site and receiving the exact inventory-and-ledger verdict.
\end{proof}

\begin{theorem}[Paper-Lean drift witness axis separation]
\label{thm:paper-lean-drift-witness-axis-separation}
For every accepted $\PaperLeanDriftWitnessUp$ carrier, an exact-resolution verdict certifies only the verification-axis lookup of the requested Lean declaration in the displayed inventory. It does not upgrade the paper-side closure axis, bridge status, or semantic theorem content attached to the paper claim.
\end{theorem}
\begin{proof}
The non-escape theorem states that the packet exports no checked theorem, statement, definition, or closure-status target beyond the exact marker-resolution rows. Its positive content is the marker-resolution verdict attached to the displayed paper marker, normalized name, label ledger, and Lean inventory lookup.

The consumer-boundary theorem confines every audit consumer to $M,A,L,I,R,H,C,P,N$. Those rows include no coordinate for changing theory closure, no bridge certificate, and no semantic equivalence row between paper claims and Lean declarations.

Thus an exact-resolution verdict may certify that the requested Lean declaration is found on the verification axis. Any paper-side closure claim, bridge status, or semantic theorem content must be supplied by a separate certificate and cannot be projected from this witness.
\end{proof}

\begin{theorem}[Paper-Lean drift witness audit-map handoff]
\label{thm:paper-lean-drift-witness-audit-map-handoff}
Every audit-map dependency consumer receiving an accepted $\PaperLeanDriftWitnessUp$ carrier receives exactly the finite packet
$$
M,A,L,I,R,H,C,P,N.
$$
It may use that packet as a row for marker-resolution dependency routing, but not as a replacement for a Lean theorem proof or for paper-side closure evidence.
\end{theorem}
\begin{proof}
The consumer boundary gives the audit-map consumer precisely the displayed finite rows $M,A,L,I,R,H,C,P,N$. These rows expose marker parsing, name normalization, duplicate-label comparison, inventory lookup, verdict exposure, transport, replay, provenance, and local naming.

Resolution stability permits the same exact-resolution outcome to be replayed under componentwise $\hsame$ transport of those rows. Closure-target resolution identifies the resolving read as the exact inventory-and-ledger verdict for the same marker site.

Axis separation then fixes the boundary of the handoff. The packet is a dependency-routing row for marker resolution; a Lean theorem proof and paper-side closure evidence remain separate obligations supplied outside the $\PaperLeanDriftWitnessUp$ witness.
\end{proof}

\begin{theorem}[Paper-Lean drift witness audit-index readiness]
\label{thm:paperleandriftwitness-audit-index-readiness}
For every accepted $\PaperLeanDriftWitnessUp$ carrier
$$
W=(M,A,L,I,R,H,C,P,N),
$$
the compact audit index for the witness is ready exactly when the same marker site exposes the exact-resolution row of $M,A,I,R$, the duplicate-label priority row of $L$, and the marker-kind separation row of $M$. The resulting index entry is a finite objectwise audit row; it is not a semantic-equivalence witness, a proof of the named Lean theorem, or evidence for the paper-side closure axis.
\end{theorem}
\begin{proof}
The exact-resolution row is supplied by \autoref{thm:paper-lean-drift-witness-resolution-exactness}: the normalized requested name in $A$ must be read against the Lean inventory $I$, and the verdict row $R$ records the exact inventory-and-ledger outcome for the same marker site in $M$.

The ledger side is supplied by \autoref{thm:paper-lean-drift-witness-duplicate-label-priority}. A refused duplicate-label row in $L$ blocks the exact export even when the inventory row contains a name-shaped hit, so the audit index cannot ignore the paper-label ledger.

The marker surface is supplied by \autoref{thm:paper-lean-drift-witness-marker-kind-separation}. A \texttt{\textbackslash leanchecked} row, a \texttt{\textbackslash leanvariant} row, a \texttt{\textbackslash leanstmt} row, a \texttt{\textbackslash leandef} row, and a \texttt{\textbackslash leantarget} row remain different paper-marker kinds.

Combining these rows gives a compact index entry over $M,A,L,I,R$. By \autoref{thm:paper-lean-drift-witness-axis-separation}, that entry remains only a verification-axis lookup row and cannot be projected into semantic equivalence, checked theorem content, or closure-axis evidence.
\end{proof}

\begin{theorem}[Paper-Lean drift witness verdict-ledger totality]
\label{thm:paperleandriftwitness-verdict-ledger-totality}
For every accepted $\PaperLeanDriftWitnessUp$ carrier, the verdict ledger over the displayed rows $M,A,L,I,R$ is total for audit consumption: each marker site is routed either through the exact inventory-and-ledger verdict, through duplicate-label refusal, or through unresolved-marker refusal. No fourth outcome is exported to a consumer.
\end{theorem}
\begin{proof}
The total verdict split is the content of \autoref{thm:paper-lean-drift-witness-inventory-verdict-totality}. It reads the displayed marker site, normalized requested name, paper-label ledger, Lean inventory lookup, and verdict row as the only accepted inputs for a marker-resolution decision.

The positive branch is constrained by \autoref{thm:paper-lean-drift-witness-resolution-exactness}; exact resolution requires the requested declaration to occur in the inventory and the same paper-label row not to be refused. This branch therefore records an objectwise inventory-and-ledger verdict, not an approximate declaration match.

The refused branches are constrained separately. Duplicate-label priority comes from \autoref{thm:paper-lean-drift-witness-duplicate-label-priority}, while absent inventory hits are confined by \autoref{thm:paper-lean-drift-witness-unresolved-marker-refusal}.

Thus every consumer-visible verdict is one of the three displayed rows. By \autoref{thm:paper-lean-drift-witness-obligation-closure-package}, the ledger remains inside the finite witness packet and exports no semantic theorem equivalence or closure-status promotion.
\end{proof}

\begin{theorem}[Paper-Lean drift witness dependency-weave handoff]
\label{thm:paperleandriftwitness-dependency-weave-handoff}
Every dependency-weave consumer receiving an accepted $\PaperLeanDriftWitnessUp$ carrier receives the finite audit packet
$$
M,A,L,I,R,H,C,P,N
$$
as a marker-resolution dependency row. The handoff preserves exact resolution, verdict-ledger totality, token objectwise audit, axis separation, and the audit-map boundary, but it does not discharge downstream semantic equivalence, Lean proof, bridge, or closure-axis obligations.
\end{theorem}
\begin{proof}
The packet delivered to the consumer is fixed by \autoref{thm:paper-lean-drift-witness-audit-map-handoff}. The dependency-weave side therefore receives marker parsing, normalized-name exposure, duplicate-label comparison, inventory lookup, verdict exposure, transport, replay, provenance, and local naming as the whole finite carrier row.

Exact resolution is preserved by \autoref{thm:paper-lean-drift-witness-closure-target-resolution}, and the verdict ledger is total by \autoref{thm:paperleandriftwitness-verdict-ledger-totality}. The handoff can route these rows through neighbouring audit-map dependencies only as objectwise marker-resolution data.

Token objectwise audit is the boundary supplied by \autoref{thm:paper-lean-drift-witness-token-objectwise-audit}: the exported row records the token object and its local audit evidence, not an ambient semantic equality between paper prose and Lean declarations.

Finally, \autoref{thm:paper-lean-drift-witness-axis-separation} and \autoref{thm:paper-lean-drift-witness-audit-map-handoff} keep the route from proving too much. Downstream semantic equivalence, Lean proof content, bridge checking, and paper-side closure status remain separate obligations outside the $\PaperLeanDriftWitnessUp$ packet.
\end{proof}

\begin{theorem}[Paper-Lean drift witness dependency-routing scope]
\label{thm:paper-lean-drift-witness-dependency-routing-scope}
Let
$$
W=(M,A,L,I,R,H,C,P,N)
$$
be an accepted $\PaperLeanDriftWitnessUp$ carrier. Every $\AuditMapDependencyWeaveUp$ consumer using $W$ as a marker-resolution dependency row must route through the displayed marker row $M$, normalized requested-name row $A$, duplicate-label ledger row $L$, Lean inventory row $I$, verdict row $R$, componentwise transport row $H$, replay row $C$, provenance row $P$, and local $\NameCert$ row $N$. The route cannot replace any of those rows by an unlisted marker, neighbouring name, outside label ledger, hidden inventory, ambient verdict, or external consumer row.
\end{theorem}
\begin{proof}
The dependency-weave handoff fixes the exported packet as $M,A,L,I,R,H,C,P,N$. It gives the audit-map consumer a dependency row for marker resolution, not a separate carrier with additional coordinates.

The audit-map handoff theorem identifies how the neighbouring consumer may read that packet. Marker parsing, name normalization, duplicate-label comparison, inventory lookup, verdict exposure, componentwise transport, replay, provenance, and local naming are the whole finite route.

The consumer boundary and non-escape rows then rule out substitutions. A consumer may transport and replay the same rows through $H,C,P,N$, but it cannot import a second marker, a different normalized name, another label ledger, an approximate inventory hit, or a verdict outside $R$.
\end{proof}

\begin{theorem}[Paper-Lean drift witness consumer-row scope]
\label{thm:paper-lean-drift-witness-consumer-row-scope}
For every accepted $\PaperLeanDriftWitnessUp$ carrier $W=(M,A,L,I,R,H,C,P,N)$, the marker, ledger, inventory, and verdict rows are grounded as finite $\BHist$ data: $M$ and $L$ are paper-site and paper-label ledger rows, $I$ is the Lean declaration inventory row, and $R$ is the exact-resolution or refusal verdict for the same marker site. Their consumer route is mediated only by componentwise $\hsame$ transport, displayed $\Cont$ replay, $\Pkg$ provenance, and the local $\NameCert$ row, and it is read by \autoref{ch:concrete-instances-auditmapdependencyweave-namecert} only as a routed dependency row.
\end{theorem}
\begin{proof}
Carrier exposure supplies the finite $\BHist$ packet. The marker row $M$, normalized-name row $A$, duplicate-label ledger $L$, inventory row $I$, verdict row $R$, transports $H$, replay routes $C$, provenance $P$, and local name $N$ are all projected from the same accepted witness.

Resolution exactness and duplicate-label priority bind the ledger and inventory reads to the same marker site. The verdict ledger totality theorem then says the visible verdict is exactly one of the accepted exact-resolution or refusal rows for $M,A,L,I$.

The only consumer movement is through $\hsame$, $\Cont$, $\Pkg$, and $\NameCert$ structure. Thus \autoref{ch:concrete-instances-auditmapdependencyweave-namecert} receives the row as routed dependency data and does not obtain a hidden theorem proof, erased obstruction, or closure-axis promotion from the witness.
\end{proof}

\begin{theorem}[Paper-Lean drift witness scoped closure package]
\label{thm:paper-lean-drift-witness-scoped-closure-package}
The $\PaperLeanDriftWitnessUp$ chapter is scoped for audit-map dependency consumption: resolution exactness fixes the inventory-and-ledger read, verdict totality exhausts the consumer-visible outcomes, axis separation prevents verification-axis rows from becoming paper-side closure evidence, and the dependency-weave handoff routes exactly the displayed packet
$$
M,A,L,I,R,H,C,P,N
$$
to the named $\AuditMapDependencyWeaveUp$ consumer surface. Within that scope, the chapter closes the paper-side dependency row for exact marker resolution, while leaving Lean verification, semantic equivalence, and bridge checking outside the witness.
\end{theorem}
\begin{proof}
First combine resolution exactness with verdict totality. The accepted witness reads one marker site through $M,A,L,I$ and exposes exactly the corresponding verdict row $R$, so the consumer-facing resolution row is objectwise fixed.

Next use axis separation. The same verdict may certify a verification-axis inventory lookup, but it cannot become a theorem proof, a semantic-equivalence row, a bridge certificate, or a paper-side closure claim outside the stated dependency scope.

Finally apply the dependency-weave handoff together with the dependency-routing and consumer-row scope theorems. The audit-map dependency consumer receives precisely $M,A,L,I,R,H,C,P,N$ through the named \autoref{ch:concrete-instances-auditmapdependencyweave-namecert} route, so the chapter supplies a scoped paper-side package for routed marker-resolution consumption and no stronger conclusion.
\end{proof}

\begin{theorem}[Paper-Lean drift witness resolution consumer readiness]
\label{thm:paper-lean-drift-witness-resolution-consumer-readiness}
For every accepted $\PaperLeanDriftWitnessUp$ carrier, any consumer route treating the witness as an exact-resolution row must read the same paper marker row $M$, normalized requested-name row $A$, paper-label ledger $L$, Lean inventory row $I$, and verdict row $R$ for one marker site. The consumer may replay those rows through $H,C,P,N$, but it cannot replace them by a neighbouring marker, a different normalized name, a different label ledger, an approximate inventory row, or a second verdict.
\end{theorem}
\begin{proof}
The exact-resolution row is fixed by \autoref{thm:paper-lean-drift-witness-resolution-exactness}. An exact verdict reads the same marker row, normalized requested name, duplicate-label ledger, Lean inventory lookup, and verdict row, and it requires the inventory hit and the absence of duplicate-label refusal on that same site.

The consumer route is bounded by \autoref{thm:paper-lean-drift-witness-consumer-boundary}. Every downstream read factors through $M,A,L,I,R,H,C,P,N$, so the consumer cannot introduce a parallel marker row or a second inventory table.

Verdict determinism from \autoref{thm:paper-lean-drift-witness-verdict-determinism} and inventory totality from \autoref{thm:paper-lean-drift-witness-inventory-verdict-totality} keep the row objectwise. Transport and replay may carry the same row through $H,C,P,N$, but the result remains the exact inventory-and-ledger verdict for the original $M,A,L,I,R$ package.
\end{proof}

\begin{theorem}[Paper-Lean drift witness duplicate-ledger readiness]
\label{thm:paper-lean-drift-witness-duplicate-ledger-readiness}
For every accepted $\PaperLeanDriftWitnessUp$ carrier, duplicate-label refusal has priority along every consumer route through $R,H,C,P,N$. If the ledger row $L$ refuses the marker site $M$, then replay, transport, provenance, and local naming may preserve that refusal, but they cannot export an exact-resolution row for that site.
\end{theorem}
\begin{proof}
Duplicate priority is the content of \autoref{thm:paper-lean-drift-witness-duplicate-label-priority}. A ledger refusal blocks exact resolution even when the normalized requested name has the shape of an inventory hit.

Unresolved-marker refusal from \autoref{thm:paper-lean-drift-witness-unresolved-marker-refusal} and non-escape from \autoref{thm:paper-lean-drift-witness-nonescape} keep refused marker sites from leaving the packet as checked markers. The route through $R,H,C,P,N$ therefore remains a refusal route whenever $L$ refuses $M$.

The dependency-weave handoff theorem, \autoref{thm:paperleandriftwitness-dependency-weave-handoff}, lets a neighbouring audit-map consumer receive the finite packet only as marker-resolution dependency data. It does not create a row that can outrank the duplicate-label ledger.
\end{proof}

\begin{theorem}[Paper-Lean drift witness inventory-refusal readiness]
\label{thm:paper-lean-drift-witness-inventory-refusal-readiness}
For every accepted $\PaperLeanDriftWitnessUp$ carrier, if the normalized requested name $A$ is absent from the Lean inventory row $I$, then every consumer-visible route through $R,H,C,P,N$ remains an unresolved-marker refusal row. The absent inventory hit cannot be transported, replayed, or repackaged as a checked marker.
\end{theorem}
\begin{proof}
Unresolved-marker refusal is supplied by \autoref{thm:paper-lean-drift-witness-unresolved-marker-refusal}. Absence of $A$ from $I$ blocks the exact-resolution condition and keeps the verdict row on the refusal side for the same marker site.

Resolution exactness from \autoref{thm:paper-lean-drift-witness-resolution-exactness} and verdict determinism from \autoref{thm:paper-lean-drift-witness-verdict-determinism} prevent a consumer from reading both an absent inventory hit and an exact-resolution verdict for the fixed rows $M,A,L,I$.

The non-escape theorem and the dependency-weave handoff theorem confine the exported material to finite marker-resolution rows. Hence replay, transport, provenance, and local naming preserve the unresolved-marker refusal rather than converting it into checked theorem content, semantic equivalence, bridge evidence, or closure-axis evidence.
\end{proof}

\begin{theorem}[Paper-Lean drift witness formalstatus sync readiness]
\label{thm:paper-lean-drift-witness-formalstatus-sync-readiness}
For every accepted $\PaperLeanDriftWitnessUp$ carrier, a downstream consumer may synchronize a $\FormalStatus$ token in the family $\formalTargetV,\encodedDefV,\scaffoldCheckedV,\theoremCheckedV,\auditCleanV,\axiomCleanV,\bridgeCheckedV$ only after reading the same paper marker row $M$, normalized requested-name row $A$, duplicate-label ledger $L$, Lean inventory row $I$, and verdict row $R$. The synchronization is an objectwise verification-axis read of the marker-resolution verdict, and it creates no paper-side closure upgrade, bridge row, semantic equivalence, or checked theorem content.
\end{theorem}
\begin{proof}
Token objectwise audit is supplied by \autoref{thm:paper-lean-drift-witness-token-objectwise-audit}. It says that the exported verification token is read from the displayed marker object and its local audit evidence, so the consumer must keep the same $M,A,L,I,R$ rows rather than replacing them by an ambient status token.

Axis separation is the boundary supplied by \autoref{thm:paper-lean-drift-witness-axis-separation}. A synchronized formalstatus row remains a verification-axis report and therefore cannot become paper-side closure, bridge evidence, semantic equivalence, or theorem proof content.

Finally, closure-target resolution from \autoref{thm:paper-lean-drift-witness-closure-target-resolution} fixes the inventory-and-ledger read that any status synchronization depends on. The consumer may record that read as a formalstatus token, but the only certified object is the same marker-resolution verdict row.
\end{proof}

\begin{theorem}[Paper-Lean drift witness marker-family resolution readiness]
\label{thm:paper-lean-drift-witness-marker-family-resolution-readiness}
The five marker families named by \texttt{\string\leanchecked}, \texttt{\string\leanvariant}, \texttt{\string\leansorryd}, \texttt{\string\leanstmt}, and \texttt{\string\leandef} remain separated through every synchronization consumer of an accepted $\PaperLeanDriftWitnessUp$ carrier. Each family resolves only through its own exact inventory-and-ledger verdict for the same $M,A,L,I,R$ rows, and no family may borrow exact resolution, duplicate-label refusal, or unresolved-marker refusal from another family.
\end{theorem}
\begin{proof}
Marker-kind separation is the first boundary. It keeps checked theorem markers, variant markers, sorry markers, statement markers, and definition markers as different paper-site rows even when their requested Lean names have similar surface shape.

Verdict-ledger totality then assigns the visible verdict for a fixed marker family and site. Apply \autoref{thm:paperleandriftwitness-verdict-ledger-totality}: the consumer sees exactly the ledger-and-inventory verdict for the selected marker row, not a family-erased verdict.

Duplicate-ledger readiness and inventory-refusal readiness handle the obstruction branches. By \autoref{thm:paper-lean-drift-witness-duplicate-ledger-readiness} and \autoref{thm:paper-lean-drift-witness-inventory-refusal-readiness}, duplicate labels and absent inventory hits remain refusals for the same marker family and cannot be repackaged as another family resolution.
\end{proof}

\begin{theorem}[Paper-Lean drift witness audit-map consumer totality]
\label{thm:paper-lean-drift-witness-audit-map-consumer-totality}
Every downstream $\AuditMapDependencyWeaveUp$ consumer of an accepted $\PaperLeanDriftWitnessUp$ carrier receives exactly one of exact resolution, duplicate-label refusal, or unresolved-marker refusal for a marker site. The verdict is transported only through $H,C,P,N$, and no ambient audit route may erase a duplicate-label or unresolved-marker obstruction.
\end{theorem}
\begin{proof}
The dependency-weave handoff is fixed by \autoref{thm:paperleandriftwitness-dependency-weave-handoff}. It sends the audit-map consumer the displayed finite packet and carries marker resolution only as dependency data.

Dependency-routing scope and consumer-row scope identify the allowed route. By \autoref{thm:paper-lean-drift-witness-dependency-routing-scope} and \autoref{thm:paper-lean-drift-witness-consumer-row-scope}, the consumer reads the same marker, name, ledger, inventory, and verdict rows, and the only transport surface is $H,C,P,N$.

The readiness rows above close the remaining consumer surface. Formalstatus synchronization from \autoref{thm:paper-lean-drift-witness-formalstatus-sync-readiness} cannot erase the obstruction, and marker-family readiness from \autoref{thm:paper-lean-drift-witness-marker-family-resolution-readiness} prevents a refused site from escaping through a different marker family.
\end{proof}

\begin{theorem}[Paper-Lean drift witness audit consumer kernel scope]
\label{thm:paperleandriftwitness-audit-consumer-kernel-scope}
Every $\AuditMapDependencyWeaveUp$ consumer of an accepted $\PaperLeanDriftWitnessUp$ carrier reads the witness only through the finite $\BHist$ rows $M,A,L,I,R,H,C,P,N$. The consumer handoff is the same marker row, normalized requested-name row, duplicate-label ledger, Lean inventory row, verdict row, componentwise $\hsame$ transport, $\Cont$ replay row, $\Pkg$ provenance row, and local $\NameCert$ naming row; it does not add an external proof object, approximate name repair, or closure-axis summary.
\end{theorem}
\begin{proof}
The audit-map consumer totality theorem, \autoref{thm:paper-lean-drift-witness-audit-map-consumer-totality}, identifies the only visible verdicts as exact resolution, duplicate-label refusal, or unresolved-marker refusal for the fixed marker site. It also confines the transport surface to the displayed downstream route.

The dependency-weave handoff theorem, \autoref{thm:paperleandriftwitness-dependency-weave-handoff}, names the consumer channel. Along that channel the finite packet is sent as dependency data, so the consumer can read $M,A,L,I,R,H,C,P,N$ and no unlisted carrier.

Formalstatus synchronization from \autoref{thm:paper-lean-drift-witness-formalstatus-sync-readiness} keeps any status read objectwise over the same marker, name, ledger, inventory, and verdict rows. Axis separation from \autoref{thm:paper-lean-drift-witness-axis-separation} then prevents this kernel-scoped handoff from becoming theorem content, bridge evidence, or a paper-side closure certificate.
\end{proof}
\leanchecked{BEDC.Derived.PaperLeanDriftWitnessUp.PaperLeanDriftWitness\_audit\_consumer\_kernel\_scope}

\begin{theorem}[Paper-Lean drift witness resolution ledger kernel scope]
\label{thm:paperleandriftwitness-resolution-ledger-kernel-scope}
For every accepted $\PaperLeanDriftWitnessUp$ carrier, exact resolution, duplicate-label refusal, and unresolved-marker refusal form a total consumer ledger over the finite rows $M,A,L,I,R$. No $\AuditMapDependencyWeaveUp$ consumer may replace that ledger by approximate matching, a neighbouring Lean declaration, a marker-family transfer, or a route-local guess.
\end{theorem}
\begin{proof}
Apply \autoref{thm:paper-lean-drift-witness-audit-map-consumer-totality}. The consumer receives exactly one of the three ledger outcomes for the marker site, and duplicate-label or unresolved-marker obstruction cannot be erased by the route.

The dependency-weave handoff theorem, \autoref{thm:paperleandriftwitness-dependency-weave-handoff}, keeps the ledger attached to the exported marker-resolution dependency row. Hence exact resolution remains a read of $M,A,L,I,R$ and not a lookup against a nearby Lean name.

The formalstatus synchronization theorem, \autoref{thm:paper-lean-drift-witness-formalstatus-sync-readiness}, gives only a token read after the same verdict row is fixed. By axis separation, \autoref{thm:paper-lean-drift-witness-axis-separation}, such a token cannot replace the total ledger or promote an obstruction into checked theorem content.
\end{proof}
\leanchecked{BEDC.Derived.PaperLeanDriftWitnessUp.PaperLeanDriftWitness\_resolution\_ledger\_kernel\_scope}

\begin{theorem}[Paper-Lean drift witness status token kernel boundary]
\label{thm:paperleandriftwitness-status-token-kernel-boundary}
For every accepted $\PaperLeanDriftWitnessUp$ carrier, formalstatus synchronization is only a verification-axis token read from the marker-resolution verdict over $M,A,L,I,R$. It cannot upgrade theoryclosure data, create a bridgestatus row, change the semantic theorem content of the paper claim, or supply an additional consumer obligation outside $H,C,P,N$.
\end{theorem}
\begin{proof}
Formalstatus synchronization readiness, \autoref{thm:paper-lean-drift-witness-formalstatus-sync-readiness}, already fixes the status token as an objectwise verification-axis read. The token is read after the same marker, normalized-name, ledger, inventory, and verdict rows are available.

Audit-map consumer totality from \autoref{thm:paper-lean-drift-witness-audit-map-consumer-totality} and the dependency-weave handoff theorem, \autoref{thm:paperleandriftwitness-dependency-weave-handoff}, identify the downstream consumer surface. The route carries the finite witness through $H,C,P,N$ rather than through a closure or bridge constructor.

Axis separation, \autoref{thm:paper-lean-drift-witness-axis-separation}, supplies the boundary. A synchronized status token reports verification-axis information only, so it cannot alter theoryclosure, bridgestatus, semantic theorem content, or the kernel-scoped consumer obligations of the witness.
\end{proof}

\begin{theorem}[Paper-Lean drift witness marker inventory kind lock]
\label{thm:paper-lean-drift-witness-marker-inventory-kind-lock}
For every accepted $\PaperLeanDriftWitnessUp$ carrier, each paper marker is resolved only through its own marker-kind row $M$, normalized requested-name row $A$, duplicate-label ledger $L$, Lean inventory row $I$, and verdict row $R$. A checked-theorem marker row may therefore not be consumed as a statement-only, encoded-definition, variant, or target marker row, and no audit-map consumer may repair the mismatch by approximate name lookup.
\end{theorem}
\begin{proof}
Marker-kind separation from \autoref{thm:paper-lean-drift-witness-marker-kind-separation} fixes the marker family before any inventory read is interpreted. The normalized requested-name row and the duplicate-label ledger are therefore attached to the marker kind already present in $M$.

Inventory verdict totality, \autoref{thm:paperleandriftwitness-verdict-ledger-totality}, supplies the only accepted outcomes over $M,A,L,I,R$. It reads exact resolution, duplicate-label refusal, or unresolved-marker refusal from those rows, rather than from a neighbouring declaration name.

Marker-family resolution readiness from \autoref{thm:paper-lean-drift-witness-marker-family-resolution-readiness} closes the consumer lock. Since each marker family has its own readiness route, a checked-theorem site cannot be retyped as another marker family or discharged by an approximate repair rule.
\end{proof}

\begin{theorem}[Paper-Lean drift witness duplicate obstruction route preservation]
\label{thm:paper-lean-drift-witness-duplicate-obstruction-route-preservation}
For every accepted $\PaperLeanDriftWitnessUp$ carrier, duplicate-label refusal remains attached to the same $M,A,L,I,R$ packet when the witness is transported through $H,C,P,N$. Audit-map dependency routing may expose the obstruction to downstream consumers, but it may not erase the duplicate ledger, substitute an exact-resolution row, or move the refusal to a different marker site.
\end{theorem}
\begin{proof}
Duplicate-label priority from \autoref{thm:paper-lean-drift-witness-duplicate-label-priority} fixes the duplicate branch before unresolved-marker lookup or exact-resolution export can be used. Duplicate-ledger readiness, \autoref{thm:paper-lean-drift-witness-duplicate-ledger-readiness}, records that branch in the same finite ledger rows $M,A,L,I,R$.

The dependency-weave handoff theorem, \autoref{thm:paperleandriftwitness-dependency-weave-handoff}, transports the witness as dependency data through $H,C,P,N$. This handoff does not introduce a second verdict row, so the duplicate obstruction travels with the original marker, name, ledger, inventory, and verdict packet.

Audit-map consumer totality, \autoref{thm:paper-lean-drift-witness-audit-map-consumer-totality}, supplies the downstream exclusion. A consumer is total over exact resolution, duplicate-label refusal, and unresolved-marker refusal, but totality does not authorize replacing the duplicate branch by a route-local exactness claim.
\end{proof}

\begin{theorem}[Paper-Lean drift witness status consumer exactness]
\label{thm:paper-lean-drift-witness-status-consumer-exactness}
For every accepted $\PaperLeanDriftWitnessUp$ carrier, a formalstatus synchronization consumer may read only the exact verdict row for the same marker site. The consumer cannot infer theoryclosure, bridgestatus, semantic theorem content, or any marker-resolution fact outside the $M,A,L,I,R,H,C,P,N$ rows of that witness.
\end{theorem}
\begin{proof}
Formalstatus synchronization readiness, \autoref{thm:paper-lean-drift-witness-formalstatus-sync-readiness}, makes the status read objectwise over the same marker, normalized requested name, duplicate ledger, inventory lookup, and verdict. The consumer therefore starts from the exact verdict row already fixed for that marker site.

The status token kernel boundary theorem, \autoref{thm:paperleandriftwitness-status-token-kernel-boundary}, keeps this read on the verification axis. It permits a status token only as a rowwise consequence of the local marker-resolution verdict, not as a constructor for closure or bridge data.

Axis separation from \autoref{thm:paper-lean-drift-witness-axis-separation} and closure-target resolution from \autoref{thm:paper-lean-drift-witness-closure-target-resolution} give the final exclusions. A synchronized formalstatus token cannot prove theoryclosure, bridgestatus, semantic theorem content, or any marker-resolution assertion beyond the exact rows supplied by the carrier.
\end{proof}

\begin{theorem}[Paper-Lean drift witness public resolution export]
\label{thm:paper-lean-drift-witness-public-resolution-export}
For every accepted $\PaperLeanDriftWitnessUp$ carrier, the public audit export of exact marker resolution consists only of the shared rows $M,A,L,I,R,H,C,P,N$. The export exposes the marker row, normalized requested-name row, duplicate-label ledger, Lean inventory row, verdict row, componentwise $\hsame$ transport, $\Cont$ replay row, $\Pkg$ provenance row, and local $\NameCert$ naming row, and it exposes no private repair rule, family-erased lookup, aggregate bridge verdict, or closure-axis certificate.
\end{theorem}
\begin{proof}
Resolution exactness from \autoref{thm:paper-lean-drift-witness-resolution-exactness} fixes the exact branch as a read of $M,A,L,I,R$. Duplicate-ledger readiness and inventory-refusal readiness, \autoref{thm:paper-lean-drift-witness-duplicate-ledger-readiness} and \autoref{thm:paper-lean-drift-witness-inventory-refusal-readiness}, show that the two obstruction branches use the same public rows rather than a hidden fallback.

The audit consumer kernel scope theorem, \autoref{thm:paperleandriftwitness-audit-consumer-kernel-scope}, identifies the only exported downstream surface as $H,C,P,N$ over the same finite witness. Hence public audit consumers receive the listed rows and no auxiliary carrier.

Finally, marker-family resolution readiness prevents a public consumer from erasing the marker family before reading the verdict. Status-token kernel boundary then prevents the exported exact-resolution row from becoming an aggregate bridge status or a paper-side closure certificate.
\end{proof}
\leanchecked{BEDC.Derived.PaperLeanDriftWitnessUp.PaperLeanDriftWitness\_public\_resolution\_export}

\begin{theorem}[Paper-Lean drift witness bridge-ready dependency surface]
\label{thm:paper-lean-drift-witness-bridge-ready-dependency-surface}
Every downstream bridge or audit consumer of an accepted $\PaperLeanDriftWitnessUp$ carrier factors through the $\AuditMapDependencyWeaveUp$ dependency surface before using the local marker-resolution verdict. Such a consumer may transport exact resolution, duplicate-label refusal, or unresolved-marker refusal along $H,C,P,N$, but it cannot replace the local verdict by an aggregate bridge status, neighbouring audit summary, or global synchronization token.
\end{theorem}
\begin{proof}
The dependency-weave handoff theorem, \autoref{thm:paperleandriftwitness-dependency-weave-handoff}, names the route through which neighbouring consumers receive the packet. Audit-map consumer totality, \autoref{thm:paper-lean-drift-witness-audit-map-consumer-totality}, then says that the received data is exactly one local verdict for the fixed marker site.

The public resolution export theorem above fixes the carrier material available to the consumer as the rows $M,A,L,I,R,H,C,P,N$. Since the dependency surface is rowwise, the consumer can use $\hsame$ transport, $\Cont$ replay, $\Pkg$ provenance, and $\NameCert$ naming only as dependency data attached to the same local verdict.

The status-token kernel boundary theorem, \autoref{thm:paperleandriftwitness-status-token-kernel-boundary}, supplies the exclusion. A synchronized token or aggregate bridge report is not a replacement verdict, so it cannot discharge duplicate-label refusal, unresolved-marker refusal, or the exactness obligation of the local marker row.
\end{proof}

\begin{theorem}[Paper-Lean drift witness public token consumer totality]
\label{thm:paper-lean-drift-witness-public-token-consumer-totality}
Every public verification-token consumer of an accepted $\PaperLeanDriftWitnessUp$ carrier is total over the three public marker-resolution outcomes: exact resolution, duplicate-label refusal, and unresolved-marker refusal. The consumer either reads the formalstatus token from the exact verdict over $M,A,L,I,R$, reads the duplicate-label refusal from the same ledger, or reads the unresolved-marker refusal from the same inventory lookup; no public token consumer is undefined on one of these branches.
\end{theorem}
\begin{proof}
Verdict-ledger totality, \autoref{thm:paperleandriftwitness-verdict-ledger-totality}, gives the three-way ledger for a fixed marker site. Resolution exactness supplies the exact branch, while duplicate-ledger readiness and inventory-refusal readiness supply the two refusal branches.

Marker-family resolution readiness keeps this totality family-local. A token consumer therefore cannot borrow an exact verdict from another marker family or reinterpret a refused branch as a different public marker kind.

Audit-map consumer totality and formalstatus synchronization readiness identify the token-facing consumer. The status-token kernel boundary theorem then keeps the token read on the verification axis, so totality over exact resolution, duplicate-label refusal, and unresolved-marker refusal does not create bridge evidence or paper-side closure data.
\end{proof}

\begin{theorem}[Paper-Lean drift witness audit replay source lock]
\label{thm:paper-lean-drift-witness-audit-replay-source-lock}
For every accepted $\PaperLeanDriftWitnessUp$ carrier, every audit replay consumed by an $\AuditMapDependencyWeaveUp$ route remains sourced in the same $M,A,L,I,R,H,C,P,N$ packet. The replay may expose the local marker row, normalized requested-name row, duplicate ledger, inventory lookup, verdict, transport, continuity replay, package provenance, and local $\NameCert$ row, but it may not substitute a neighbouring marker, neighbouring normalized name, or neighbouring verdict for the original packet.
\end{theorem}
\begin{proof}
Resolution exactness, \autoref{thm:paper-lean-drift-witness-resolution-exactness}, first fixes the exact branch as a rowwise read of $M,A,L,I,R$. Marker inventory kind lock, \autoref{thm:paper-lean-drift-witness-marker-inventory-kind-lock}, keeps the marker family and inventory lookup attached to that same source before any consumer sees the replay.

The public resolution export theorem, \autoref{thm:paper-lean-drift-witness-public-resolution-export}, names the exported rows as $M,A,L,I,R,H,C,P,N$. Applying public token consumer totality then shows that exact resolution, duplicate-label refusal, and unresolved-marker refusal are all consumed from this single packet, not from a second witness.

Finally, the dependency-weave handoff route only transports the exported rows to neighbouring audit-map consumers. Since the handoff receives the locked packet after marker-kind separation and exactness, it cannot replace the source marker or verdict by a neighbouring marker summary.
\end{proof}

\begin{theorem}[Paper-Lean drift witness Lean inventory hit exhaustion]
\label{thm:paper-lean-drift-witness-lean-inventory-hit-exhaustion}
For every accepted $\PaperLeanDriftWitnessUp$ carrier, exact marker resolution is exhausted by the displayed Lean inventory hit together with the non-duplicate label row. If the duplicate-label ledger marks the requested name as duplicated, or if the Lean inventory lookup is absent, the corresponding branch remains a public refusal row and cannot be consumed as an exact hit.
\end{theorem}
\begin{proof}
Resolution exactness identifies the exact branch with the normalized requested-name row $A$, duplicate-label ledger $L$, Lean inventory lookup $I$, and verdict row $R$. The exact branch is therefore available only when the inventory hit and the non-duplicate label condition are both present in the displayed packet.

Marker inventory kind lock keeps this inventory read within the marker family fixed by $M$. Public resolution export then exposes only the same inventory and ledger rows to downstream consumers, so no approximate declaration name or marker-family coercion can supply a second exact hit.

Public token consumer totality supplies the exhaustive case split. The exact branch consumes the displayed hit; the duplicate-label branch consumes the duplicate ledger; the unresolved-marker branch consumes the absent inventory lookup. The two refusal branches remain refusal rows after export and cannot be reclassified as exact resolution.
\end{proof}

\begin{theorem}[Paper-Lean drift witness name normalization no-escape]
\label{thm:paper-lean-drift-witness-name-normalization-noescape}
For every accepted $\PaperLeanDriftWitnessUp$ carrier, underscore and Lean-name normalization is only a replay row inside the same $M,A,L,I,R,H,C,P,N$ packet. Normalization authorizes comparison with the displayed requested-name row, but it gives no permission for approximate lookup, spelling repair, marker-family coercion, or substitution of a neighbouring declaration name.
\end{theorem}
\begin{proof}
Resolution exactness uses normalization as part of the requested-name row $A$ and then reads the duplicate ledger, inventory lookup, and verdict from $L,I,R$. Thus normalization is an input row to exact resolution, not an external search rule.

Marker inventory kind lock keeps the normalized name under the marker family fixed by $M$. Public resolution export carries this normalized name together with $H,C,P,N$, so audit-map consumers receive the replayed spelling only as part of the same finite packet.

Public token consumer totality closes the exclusion. A consumer is total over exact resolution, duplicate-label refusal, and unresolved-marker refusal for the given packet, and this totality leaves no branch in which normalization can escape into approximate lookup or coerce one marker family into another.
\end{proof}

The exact-resolution row exported by $\PaperLeanDriftWitnessUp$ is consumed by the cross-map routing and synthesis-consumer rows of \autoref{ch:concrete-instances-auditmapdependencyweave-namecert}. The $\PaperLeanDriftWitnessUp$ packet supplies marker normalization, duplicate-label refusal, Lean inventory lookup, and exact-resolution verdict rows, while $\AuditMapDependencyWeaveUp$ decides how such local rows travel through neighbouring audit-map routes without discharging their obstruction rows.

\closureat{PaperLeanDriftWitnessUp}{seedStr}
\closureat{PaperLeanDriftWitnessUp}{obligationStr}
\closureat{PaperLeanDriftWitnessUp}{scopedStr}

\begin{closurestatus}{\PaperLeanDriftWitnessUp}
  \constructivestory{A $\PaperLeanDriftWitnessUp$ packet is generated from a paper marker row, normalized requested-name row, duplicate-label ledger, Lean inventory lookup, resolution verdict, componentwise $\hsame$ transport, displayed replay routes, $\Pkg$ provenance, and a local $\NameCert$ row.}
  \theoryclosure{\scopedClosure}
  \scopeclosed{The scoped surface binds the exact-resolution, duplicate-label, unresolved-marker, marker-kind, inventory-verdict, transport, replay, provenance, dependency-weave handoff, and audit-map dependency consumer rows for exact paper-Lean marker resolution.}
  \formalstatus{\unformalizedV}
  \bridgestatus{none}
  \origin{human}
  \notclaimed{This chapter does not close semantic equivalence between paper claims and Lean declarations, checked theorem content, bridge checking, or any permission to treat unresolved markers as verified.}
  \upgradepath{Public closure requires Lean verification of the scoped dependency package and a public bridge surface for downstream audit consumers.}
\end{closurestatus}
